\documentclass[11pt,twoside]{article}

\usepackage[
  a4paper,
  margin=1.5cm,
  includehead, includefoot
]{geometry}

\usepackage[utf8]{inputenc}
\usepackage[brazil,portuges]{babel}
\usepackage{csquotes}

% \usepackage[
%   backend=biber,
%   style=abnt,
%   language=auto,
% ]{biblatex}

%   \addbibresource{../refs/referencias.bib}

\usepackage{soul}
\usepackage{xcolor}
\usepackage{framed}

\usepackage{tikz}

  \usetikzlibrary{shapes.geometric, positioning, decorations.pathreplacing}

\usepackage{kpfonts}
\usepackage[T1]{fontenc}

\usepackage{tabularx}

  \newcolumntype{C}{>{\centering\arraybackslash}X}

\usepackage{graphicx}
\usepackage{multirow}
\usepackage{multicol}
\usepackage{hhline}
\usepackage{booktabs}
\usepackage[inline]{enumitem}
\usepackage{xfrac}
\usepackage{mathtools}

  \newcommand{\Binary}{\ensuremath{\mathbb{B}}}
  \newcommand{\Natural}{\ensuremath{\mathbb{N}}}
  \newcommand{\Integer}{\ensuremath{\mathbb{Z}}}
  \newcommand{\Rational}{\ensuremath{\mathbb{Q}}}
  \newcommand{\Real}{\ensuremath{\mathbb{R}}}

  \DeclarePairedDelimiter{\floor}{\lfloor}{\rfloor}
  \DeclarePairedDelimiter{\ceil}{\lceil}{\rceil}
  \DeclarePairedDelimiter{\abs}{\lvert}{\rvert}

\usepackage{amsthm}

  \newtheorem{proposition}{Proposição}[section]
  \theoremstyle{definition}
  \newtheorem{definition}{Definição}[section]

\usepackage{multicol}

  \newcommand{\finishcols}{\vfill\vphantom{.}}

\usepackage[portuguese,algoruled,vlined,noend]{algorithm2e}

  \SetKwInput{Algoritmo}{Algoritmo}
  \SetKw{Retorne}{retorne}
  \SetKwFor{Para}{para}{faça}{fimpara}
  \SetKw{Escreva}{escreva}
  \DontPrintSemicolon

\usepackage{hyperref}
\usepackage[brazil]{varioref}
\usepackage[brazilian]{cleveref}

  \def\disciplinacodigotext{QXD0115}
  \def\disciplinanometext{Estrutura de Dados Avançada}
  \def\semestretext{2022.1}
  \def\provatype{Notas de Aula}
  \def\provatitle{Rotações e Balanceamento de AVLs}
  \def\provatext{\provatype{} \ifx\provatype\empty {} \else {--} \fi \provatitle{}}
  \def\provacodtext{}
  \def\professorabbrv{Prof.}
  \def\professornometext{Arthur Araruna}

\usepackage{fancyhdr}

  \pagestyle{fancy}
  \fancyhead{}
  \lhead{\disciplinacodigotext}
  \chead{\provatext}
  \rhead{\semestretext}
  \fancyfoot{}
  \fancyfoot[LE,RO]{\colorbox{black!20}{\thepage}}

% * NOTE: Redefinição feita para evitar confusão com exponenciações.
\renewcommand{\thefootnote}{(\roman{footnote})}

\begin{document}
  \noindent
  \begin{minipage}{\textwidth}
    \begin{tabularx}{\textwidth}{cXc}
      \multirow{4}{*}{\includegraphics[height=1.8cm]{brasaoufc.pdf}} & {\bf UNIVERSIDADE FEDERAL DO CEARÁ}         & \multirow{3}{*}{} \\
                                                                     & {\bf Campus de Quixadá}                     &                   \\
                                                                     & \professorabbrv\enspace\professornometext   &                   \\
                                                                     & \disciplinacodigotext - \disciplinanometext & {}                \\
    \end{tabularx}
  \end{minipage}

  \part*{\provatitle}

    \section{Conceitos}

     Vamos relembrar o nossos conceitos que serão necessários para o resultado a ser demonstrado:

     \begin{definition}

     \end{definition}

    \section{Conclusão}


     %\part*{Apêndice}


     %\printbibliography

\end{document}